\documentclass[12pt]{article}
\usepackage[margin=0.3in,top=0.25in,bottom=0.4in]{geometry}
\makeatletter
\renewcommand{\@maketitle}{\newpage\null\vskip -3.5em\begin{center}{\LARGE \@title \par}\end{center}\par\vskip 0.2em}
\makeatother
\usepackage{amsmath,amssymb}
\usepackage{booktabs}
\usepackage{hyperref}
\setlength{\parskip}{0.6em}
\setlength{\parindent}{0pt}

\title{Why our scale uncertainty is small at low $p_{D\perp}$}
\author{}
\date{}

\begin{document}
\maketitle

\section*{In short}

In our $D^0$ predictions the scale-variation band is narrowest at low transverse momentum,
while in light-hadron calculations such as Casuga, Mulani and M\"antysaari
(arXiv:2609.15461) it is widest there. This is not a bug. It comes from three features of our
calculation: the scale is never small, the fragmentation scale cannot go below the charm mass,
and the exclusive part, which is large at low $p_{D\perp}$, does not depend on $\alpha_s$. As a
consequence, our band underestimates the theory uncertainty at low $p_{D\perp}$, and the paper
should say so.

\section*{1. What the light-hadron calculation does}

In arXiv:2609.15461 (charged hadrons in DIS), the fragmentation scale is the hadron transverse
momentum, $\mu=p$, varied between $p/2$ and $2p$, with the light-quark NNFF1.0 functions. At low
$p$ the scale is very low, around $0.5$--$1$~GeV. There $\alpha_s$ is large and the fragmentation
functions change quickly with $\mu$, so a factor 2 in $\mu$ gives a large change: the band grows
at low $p$.

\section*{2. Our scale is never small}

We use $\mu=m_t=\sqrt{p_{D\perp}^2+m^2}$ with $m=1.5$~GeV. Even at $p_{D\perp}=0$ the scale is
$1.5$~GeV, and near this scale the charm fragmentation functions have hardly evolved from their
starting point $\mu_0=m$. Their dependence on $\mu$ is therefore weak at low $p_{D\perp}$.

\section*{3. Half of the $\mu_F$ variation is cut off at low $p_{D\perp}$}

The charm fragmentation functions start at the charm mass (BCFY and KK at $\mu_0=m=1.5$~GeV;
HymnD is zero below its $m_c=1.3$~GeV). The code therefore does not allow a fragmentation scale
below $m$ (\texttt{src/main.cpp}: \texttt{if (frag\_scale < param.m) frag\_scale = param.m}).
For the lower variation $\mu_F=\mu/2$ this happens whenever $\mu/2<m$, i.e.\ for
$p_{D\perp}<\sqrt3\,m\approx2.6$~GeV:

\begin{center}
\begin{tabular}{cccc}
\toprule
$p_{D\perp}$ [GeV] & $\mu=m_t$ [GeV] & $\mu/2$ [GeV] & $\mu_F$ actually used for ``$\mu/2$'' [GeV] \\
\midrule
0   & 1.50 & 0.75 & 1.50 (cut) \\
1   & 1.80 & 0.90 & 1.50 (cut) \\
2   & 2.50 & 1.25 & 1.50 (cut) \\
2.6 & 3.00 & 1.50 & 1.50 \\
5   & 5.22 & 2.61 & 2.61 \\
10  & 10.11 & 5.06 & 5.06 \\
\bottomrule
\end{tabular}
\end{center}

So in the low-$p_{D\perp}$ bins, ``$\mu_F=\mu/2$'' is the same as $\mu_F=m$, close to the central
value. Only the upward variation $\mu_F=2\mu$ changes the result there. The cut is physical,
since the charm fragmentation functions are not defined below the charm mass, but it halves the
variation at low $p_{D\perp}$.

\section*{4. $\mu_R$ only affects part of the cross section}

$\alpha_s(\mu_R)$ only multiplies the diffractive (TMD) contribution. The exclusive contribution
has no $\alpha_s$: the strength of the strong interaction is contained in the dipole amplitude,
which is fitted to HERA data. At low transverse momentum the exclusive part is large (for the
charm quark it dominates below $p_\perp\approx5$~GeV), so varying $\mu_R$ moves only a fraction of
the total cross section there.

\section*{5. Suggested text for the paper}

\begin{quote}
Since the charm fragmentation functions are defined only above the charm-quark mass, the
factorization scale is not allowed to go below $\mu_0=m$; for $p_{D\perp}\lesssim\sqrt3\,m\approx
2.6$~GeV the lower variation $\mu_F=\mu/2$ is therefore replaced by $\mu_F=m$. Moreover,
$\alpha_s(\mu_R)$ enters only the diffractive TMD contribution, while the exclusive one, which
dominates at low $p_{D\perp}$, does not depend on $\mu_R$. As a consequence, the scale uncertainty
is smallest at low $p_{D\perp}$ and should be regarded as a lower estimate there.
\end{quote}

\section*{6. Options for a more realistic low-$p_{D\perp}$ uncertainty}

\begin{itemize}
\item \textbf{Fragmentation-model band:} the spread between BCFY, KK and HymnD. It is larger than
  the scale band at all $p_{D\perp}$ (BCFY is $1.4$--$1.6$ times KK and HymnD after the
  normalization fix) and needs no new runs.
\item \textbf{HymnD replica band:} the HymnD replicas are widest at $z_h\lesssim0.5$ near the
  initial scale, which matters most at low $p_{D\perp}$. The replica runs exist on Roihu (the
  \texttt{hymnd\_band} step of the submit script).
\end{itemize}

\end{document}
